\documentclass[12pt]{article}
\usepackage{latexsym}
\usepackage{amssymb,amsmath}
\usepackage[pdftex]{graphicx}
\usepackage{tikz}

\topmargin = 0.1in \textwidth=5.7in \textheight=8.6in

\oddsidemargin = 0.2in \evensidemargin = 0.2in

\begin{document}

\begin{center}
    COMPUTER SCIENCE E-20, SPRING 2014 \\
    Homework Problems\\
    Strong Induction, Induction Review, Propositional Logic \\
    Author: Tawheed Abdul-Raheem
\end{center}

\smallskip

\begin{enumerate}

    \item Suppose you are given a real number $x$ such that $x + \frac{1}{x}$ is an integer.  Use Strong Induction to show that $x^n + \frac{1}{x^n}$ is an integer for all integers $n$.

        \textbf{Proof: }
        We will prove by strong induction that $x^n + \frac{1}{x^n}$ is an integer for all integers $n$

        \textbf{Base case: } Lets assume that $(n = 0): P(0)$
        \[ x^0 + \frac{1}{x^0}\]
        The solution we get from above brings us back to our proposition where $x$ is a real number

        \textbf{Basis step: } Lets assume that $(n = 1): P(1)$
        \[ x^1 + \frac{1}{x^1}\]

        \textbf{Inductive step: } Suppose that $n \geq 1$ represents integers that can fit in $x + \frac{1}{x}$ such $x$ is a real number. We must show that $P(n+1)$
        also holds for this argument, namely $n + 1$ otherwise the argument does not hold

        \textbf{Conclusion: } Since adding $1$ to $n$ makes the number positive its fair to say that our argument holds for when $P(n) is (n+1)$

    \item Prove using strong induction that any square can be subdivided into $n$ smaller squares, where $n>5$. For example, the large square below has been subdivided into 6 squares.\\
        {\em Hint: first show that any square subdivided into $k$ squares can easily be subdivided into $k+3$ squares, then think how many base cases you need show are true (it is not just the case of n=6).}

        \begin{center}
            \begin{tikzpicture}

                \draw (0,0) -- (3,0) -- (3,3) -- (0,3) -- (0,0);
                \draw (2,0) -- (2,2) -- (0,2);
                \draw (1,2) -- (1,3);
                \draw (2,2) -- (2,3);
                \draw (2,1) -- (3,1);
                \draw (3,2) -- (2,2);

            \end{tikzpicture}
        \end{center}
        \textbf{Solution: } Each of the squares can be easily subdivided to $n+3$. When we we subdivided the square, we get additional $3$ squares, adding that to the original square we have $n+3$ sub squares. Now we need to show that our squares can be subdivided in to the $7$ and $8$. We can subdivide our $7$ and $8$ easily with $n+3$

    \item The Fibonacci numbers are defined by $F_0 = 0, F_1 = 1, F_n = F_{n-1} + F_{n-2}$.  Prove  using strong induction that for all $n \geq 2$:
        \begin{equation*}
            F_n \geq \phi^{n-2}
        \end{equation*}
        where $\phi$ is the constant $\frac{1+\sqrt{5}}{2}$.  \\
        {\it Hint:} Verify and use the fact that $\phi + 1 = \phi^2$.

        \textbf{Solution: } We will use mathematical induction to prove that for all integers $n \geq 2$,$P(n)$ is true
        So first we show that $P(0)$ and $P(1)$ are true we let $P(0)$ be $F_0$ and $P(1)$ be $F_1$ \\

        \textbf{ Lets prove P(2) }
        \[ F_2  \geq \phi^{n-2} \]
        \[ F_2  \geq \phi^{0} \]
        \[ 1  \geq 1 \]

        \textbf{ Lets prove P(3) }
        \[ F_3  \geq \phi^{n-2} \]
        \[ F_3  \geq \phi^{3-2} \]
        \[ 2 \geq \frac{1+\sqrt{5}}{2} \]

        \textit{\textbf{
                for all integer $n \geq 2$, if $P(n)$ is true for all integers 
                n from 2 then it must be true for $P(n+1)$
            }}
        Let this be our predicate that we are trying to prove for $n+1$
        \[ F_{n+1} = \phi^{n-1} \]
        The above was derived from substituition.

        \[ F_{n} \geq \frac{\phi^{(n)}}{ \phi + 1} \]
        \[ F_{n-1} \geq \frac{\phi^{(n-1)}}{ \phi + 1} \]
        \textbf{\it{ By summing those two inequalities we get }}
        \[ F_{n} + F_{n-1} \geq \frac{\phi^{(n)}}{ \phi + 1} + \frac{\phi^{(n-1)}}{ \phi + 1} \]
        \[ F_{n+1} \geq \frac{(\phi^{n-1})(\phi + 1)}{ \phi + 1} \]
        \[ F_{n+1} \geq \phi^{n-1} \]

       We have successfully proved our theorem.
       

    \item
        \begin{enumerate}

            \item Define the propositions p=``You obey the speed limit'' and the q=``You are going to a wedding''. Write the following sentences as compound propositions using p and q: 
                \begin{enumerate}
                    \item Failing to obey the speed limit implies that you are going to a wedding.
                    \item You drive below the speed limit only if you are going to a wedding. 
                    \item You do not obey the speed limit unless you are going to a wedding. 
                    \item You drive above the speed limit whenever you are going to a wedding. 
                \end{enumerate}

                \textbf{Solution: } \\
                p=``You obey the speed limit'' \\
                q=``You are going to a wedding''
                \begin{enumerate}
                    \item $\neg p \rightarrow q $
                    \item $p \leftrightarrow q$
                    \item $p \leftrightarrow q$
                    \item $q \rightarrow \neg p$
                \end{enumerate}

            \item Define the propositions p=``The home team wins,'' q=``It is raining,'' r=``There is an earthquake '' Write the following sentences as compound propositions using p,q, and r:
                \begin{enumerate}
                    \item Rain and earthquake are sufficient for the home team to win. 
                    \item Rain and earthquake are necessary but not sufficient for the home team to win. 
                    \item The home team wins only if it is not raining and there is no earthquake.
                    \item If it is raining the home team will win unless there is an earthquake.  
                \end{enumerate}

                \textbf{Solution: } \\
                p=``The home team wins,'' \\
                q=``It is raining,'' \\
                r=``There is an earthquake ''
                \begin{enumerate}
                    \item $(q \wedge r) \rightarrow p $
                    \item $p \rightarrow (q \wedge r) $
                    \item $p \rightarrow \neg(q \wedge r) $
                    \item $ (q \wedge \neg r) \rightarrow p$
                \end{enumerate}

        \end{enumerate}

    \item Using a truth table, determine which of the following are equivalent to $(p \wedge q) \rightarrow r$ and which are equivalent to $(p \vee q) \rightarrow r$:
\begin{enumerate}
    \item $p \rightarrow (q \rightarrow r)$
    \item $q \rightarrow (p \rightarrow r) $
    \item $(p \rightarrow r) \wedge (q \rightarrow r) $
    \item $(p \rightarrow r) \vee (q \rightarrow r)$
\end{enumerate}

\textbf{Solution: } \\
\begin{displaymath}
\begin{array}{|c|c|c|c|c|c|c|}
    p & q & r & p \wedge q & p \vee q & (p \wedge q) \rightarrow r & (p \vee q) \rightarrow r \\
\hline
T & T & T & T & T & T & T \\
F & T & T & F & T & T & T \\
T & F & T & F & T & T & T \\
F & F & T & F & F & T & T \\
T & T & F & T & T & F & F \\
F & T & F & F & T & T & F \\
T & F & F & F & T & T & F \\
F & F & F & F & F & T & T \\
\end{array}
\end{displaymath}


\begin{displaymath}
\begin{array}{|c|c|c|c|c|c|c|}
     q \rightarrow r &  p \rightarrow (q \rightarrow r) & p \rightarrow r & q \rightarrow (p \rightarrow r) & (p \rightarrow r) \wedge (q \rightarrow r) & (p \rightarrow r) \vee (q \rightarrow r) \\
\hline
T & T & T & T & T & T \\
T & T & T & T & T & T \\
T & T & T & T & T & T \\
T & T & T & T & T & T \\
F & F & F & F & F & F \\
F & T & T & T & F & T \\
T & T & F & T & F & T \\
T & T & T & T & T & T \\
\end{array}
\end{displaymath}

Equivalent to $(p \wedge q) \rightarrow r$
\begin{enumerate}
    \item $p \rightarrow (q \rightarrow r)$
    \item $q \rightarrow (p \rightarrow r) $
    \item $(p \rightarrow r) \vee (q \rightarrow r)$
\end{enumerate}


Equivalent to $(p \vee q) \rightarrow r$
\begin{enumerate}
    \item $(p \rightarrow r) \wedge (q \rightarrow r) $
\end{enumerate}

\end{enumerate}

\pagebreak

\end{document}
